\hypertarget{capacity_cycles-a-matrix-capacity-obstruction-and-explicit-five-phase-fermat-cycles}{%
\chapter{A matrix-capacity obstruction and explicit five-phase Fermat
cycles}\label{capacity_cycles-a-matrix-capacity-obstruction-and-explicit-five-phase-fermat-cycles}}

Komorebi research continuation, version 0.2.0. Session: 21 September
2026, America/Chicago.

\hypertarget{capacity_cycles-status-and-scope}{%
\section{Status and scope}\label{capacity_cycles-status-and-scope}}

This note contains manuscript proofs and replayable exact algebraic
certificates. It does not prove the general Hodge conjecture. The exact
certificates were independently replayed in October. The geometric
arguments are not formally verified or peer reviewed.

The positive algebraicity family reconstructed here is the \textbf{known
five-phase standard family of Aoki}. The October original-source audit
checked Theorem 2-1 of the 1987 paper. Its equations give the same
family after the Newton identities below. The construction is a direct
recomputation of those known surfaces. The determinant and period are
derived explicitly using the checked complete-intersection theorem of
Villaflor Loyola. No priority for the surfaces, equations, algebraicity
or determinant is claimed.

The concrete advances within this research attempt are:

\begin{enumerate}
\def\labelenumi{\arabic{enumi}.}
\tightlist
\item
  An exact global upper bound on the normalized weight sum of every
  invariant-column matrix candidate, with equality for the earlier
  degree-70 carrier.
\item
  A uniform obstruction showing that the earlier positive-anticanonical
  matrix criterion cannot cover all balanced Fermat packets, even with
  unbounded exponents and coordinate-power lifts.
\item
  Explicit complete-intersection surfaces, a uniform determinant
  identity, and nonzero exact periods for the infinite standard family
  missed by that criterion.
\end{enumerate}

The original v0.1.0 archive is preserved byte for byte in
\texttt{predecessor/}. Its original 24-check verifier was replayed
successfully before this continuation. The geometric arguments of the
predecessor have not been independently formally verified.

\hypertarget{capacity_cycles-conventions}{%
\section{Conventions}\label{capacity_cycles-conventions}}

Let

\[
X_m=\{F_m=x_0^m+\cdots+x_5^m=0\}\subset\mathbf P^5_{\mathbf C},
\qquad G_m=(\mu_m)^6/\mu_m.
\]

Write \(\langle r\rangle_m\) for the least-positive residue of a nonzero
residue class. A full-support character is

\[
a=(a_0,\ldots,a_5),\qquad 0<a_i<m,\qquad \sum a_i\equiv0\pmod m.
\]

The Griffiths residue basis of the Fermat hypersurface diagonalizes the
action of \(G_m\). The monomial \(x_0^{a_0-1}\cdots x_5^{a_5-1}\) at
pole order three gives the character of middle Hodge type when
\(\sum a_i=3m\). The entire rational orbit has type \((2,2)\) when

\begin{equation}\label{eq:capacity_cycles-1.1}
\sum_i\langle ta_i\rangle_m=3m
\quad\text{for every }t\in(\mathbf Z/m\mathbf Z)^\times.\end{equation}

A rational packet is the rational subrepresentation whose
complexification is the sum of these distinct conjugate eigenlines. An
individual complex eigenvector is not being identified with an
individual rational class.

The Jacobian ideal and socle monomial are

\begin{equation}\label{eq:capacity_cycles-1.2}
J_{F_m}=(x_0^{m-1},\ldots,x_5^{m-1}),\qquad
S_m=\prod_{i=0}^5x_i^{m-2}.\end{equation}

We use

\[
\Omega=\sum_{i=0}^5(-1)^ix_i\,dx_0\wedge\cdots\widehat{dx_i}\cdots\wedge dx_5
\]

and the residue normalization of Villaflor Loyola, Theorems 1.1 and 1.2.
The order of polynomial rows in each Jacobian below is explicitly fixed.

\hypertarget{capacity_cycles-the-matrix-capacity-and-a-certificate-over-an-infinite-set}{%
\section{The matrix capacity and a certificate over an infinite
set}\label{capacity_cycles-the-matrix-capacity-and-a-certificate-over-an-infinite-set}}

Fix a character \(a\) modulo \(m\). Its invariant exponent columns form

\[
\mathcal L_a=\{v\in\mathbf Z_{\ge0}^6\setminus\{0\}:a\cdot v\equiv0\pmod m\}.
\]

Define

\begin{equation}\label{eq:capacity_cycles-2.1}
C(a)=\sup\left\{\sum_vw_v:
\sum_vw_vv=\mathbf1,\quad w_v\ge0,
\quad\text{only finitely many }w_v\ne0\right\}.\end{equation}

Every nonsingular nonnegative integral matrix \(A\) satisfying
\(A^Ta\in m\mathbf Z^6\), with \(w=A^{-1}\mathbf1>0\), gives one
feasible point. Hence

\begin{equation}\label{eq:capacity_cycles-2.2}
\mathbf1^TA^{-1}\mathbf1\le C(a).\end{equation}

This is a necessary numerical bound, not a sufficient geometric
criterion. A capacity maximizer can use fewer than six columns. A
six-column maximizer can fail normality or well-formedness. No claim
about a normal carrier follows merely from \(C(a)>1\).

\hypertarget{capacity_cycles-lemma-2.1-global-dual-certificate}{%
\subsection{Lemma 2.1: global dual
certificate}\label{capacity_cycles-lemma-2.1-global-dual-certificate}}

Suppose \(y\in\mathbf Q_{>0}^6\) and \(y\cdot v\ge1\) for every
\(v\in\mathcal L_a\). Then

\[
C(a)\le y\cdot\mathbf1.
\]

Indeed,

\[
\sum_vw_v\le\sum_vw_v(y\cdot v)
=y\cdot\sum_vw_vv=y\cdot\mathbf1.
\]

The infinitely many inequalities can be certified by a finite residue
graph. Write \(y_i=c_i/L\), where \(c_i,L\) are positive integers. On
\(\mathbf Z/m\mathbf Z\), an edge \(r\to r+a_i\) has cost \(c_i\). An
integer function \(d\) satisfying

\begin{equation}\label{eq:capacity_cycles-2.3}
d(0)=0,\qquad d(r+a_i)\le d(r)+c_i,
\qquad d(-a_i)+c_i\ge L\end{equation}

for every \(r,i\) certifies the dual inequality for \textbf{every}
invariant column, regardless of its total degree.

To prove this, order the multiset described by a nonzero invariant
column into a nonempty zero-sum word. Repeated application of the first
inequalities bounds the cost of its prefix below by \(d(-a_i)\), where
\(i\) is its last letter. Adding the last cost gives at least \(L\).
Divide by \(L\).

The program uses a shortest-path algorithm to find a potential, but
verifies \eqref{eq:capacity_cycles-2.3} directly. The correctness of the global conclusion
therefore depends on the displayed finite inequalities and this proof,
not on trusting a shortest-path implementation or a search-degree
cutoff.

\hypertarget{capacity_cycles-proposition-2.2-the-old-degree-70-carrier-is-globally-optimal-in-normalized-weight-sum}{%
\subsection{Proposition 2.2: the old degree-70 carrier is globally
optimal in normalized weight
sum}\label{capacity_cycles-proposition-2.2-the-old-degree-70-carrier-is-globally-optimal-in-normalized-weight-sum}}

For

\[
a=(1,20,24,42,61,62)\pmod{70},
\]

take

\begin{equation}\label{eq:capacity_cycles-2.4}
y=\frac1{70}(11,10,19,14,6,17).\end{equation}

The exact potential in \texttt{results/capacity\_certificates.json}
verifies all 420 edge inequalities in \eqref{eq:capacity_cycles-2.3}, and all six last-edge lower
bounds are exactly 70. Thus

\[
C(a)\le\sum_i y_i=\frac{11}{10}.
\]

The matrix from v0.1.0 is

\[
A=\begin{pmatrix}
0&0&2&0&0&3\\
0&0&1&2&7&0\\
1&0&2&2&0&1\\
0&5&0&0&0&0\\
0&0&0&2&0&3\\
3&0&0&0&0&0
\end{pmatrix},
\quad
A^{-1}\mathbf1=(1/3,1/5,1/10,1/10,1/10,4/15).
\]

It is feasible and its weights sum to \(11/10\). Therefore

\begin{equation}\label{eq:capacity_cycles-2.5}
\boxed{C(1,20,24,42,61,62\bmod70)=11/10.}\end{equation}

This is global within the invariant-column matrix model, not merely
optimal among columns of degree at most 18. It is not an assertion that
a scale-dependent anticanonical index has an invariant maximum.

\begin{figure}[tbp]\centering\includegraphics[width=0.98\linewidth]{26_capacity_values.pdf}\caption{Seven named capacity certificates and the degree-210 target. These eight specified characters are not a census of their degrees.}\end{figure}

\hypertarget{capacity_cycles-a-uniform-obstruction-to-universal-matrix-carrier-coverage}{%
\section{A uniform obstruction to universal matrix-carrier
coverage}\label{capacity_cycles-a-uniform-obstruction-to-universal-matrix-carrier-coverage}}

For \(q\ge2\) and \(1\le b<q\), set

\begin{equation}\label{eq:capacity_cycles-3.1}
m=5q,\qquad
 a_{q,b}=(b,b+q,b+2q,b+3q,b+4q,5(q-b)).\end{equation}

\hypertarget{capacity_cycles-lemma-3.1-balance}{%
\subsection{Lemma 3.1: balance}\label{capacity_cycles-lemma-3.1-balance}}

Every character in \eqref{eq:capacity_cycles-3.1} is full support and its rational orbit is of
type \((2,2)\).

Let \(t\) be a unit modulo \(5q\), and let \(b'=tb\bmod q\), with
\(1\le b'<q\). Multiplication by \(t\) permutes the five residues in the
fiber over \(b'\), because \(t\) is a unit modulo 5. The first five
least-positive residues are therefore, in some order,

\[
b',b'+q,b'+2q,b'+3q,b'+4q.
\]

The last residue is \(5(q-b')\). Their sum is \(15q=3m\). This is a
proof for all parameters, separate from the finite checks in the bundle.

\hypertarget{capacity_cycles-no-positive-anticanonical-invariant-column-matrix}{%
\subsection{No positive-anticanonical invariant-column
matrix}\label{capacity_cycles-no-positive-anticanonical-invariant-column-matrix}}

Assume \(q>5\) and \(\gcd(b,q)=1\). For every nonsingular nonnegative
integral matrix \(A\) such that

\[
A^Ta_{q,b}\in5q\mathbf Z^6,
\qquad w=A^{-1}\mathbf1>0,
\]

one has

\begin{equation}\label{eq:capacity_cycles-3.2}
\boxed{\sum_j w_j\le1.}\end{equation}

In fact \(C(a_{q,b})=1\).

\textbf{Proof.} Define

\[
y_q=\frac1q(1,1,1,1,1,q-5),\qquad y_q\cdot\mathbf1=1.
\]

If \(v\in\mathcal L_{a_{q,b}}\), reduction modulo \(q\), followed by
cancellation of the unit \(b\), gives

\[
v_0+v_1+v_2+v_3+v_4-5v_5\equiv0\pmod q.
\]

Hence

\[
N=v_0+v_1+v_2+v_3+v_4+(q-5)v_5
\]

is divisible by \(q\). It is a positive integer because \(v\ne0\) and
\(q>5\). Therefore \(N\ge q\), or \(y_q\cdot v\ge1\). Lemma 2.1 proves
\(C(a_{q,b})\le1\).

Conversely \(v=\mathbf1\) is invariant, because \(\sum a_i=3m\), and
putting weight 1 on it is a feasible point of \eqref{eq:capacity_cycles-2.1}. Thus
\(C(a_{q,b})=1\), proving \eqref{eq:capacity_cycles-3.2}. \(\square\)

The invariant column set is unchanged by multiplying a character by a
unit modulo \(m\), or by replacing \((m,a)\) by \((km,ka)\). A
coordinate permutation merely permutes the columns' entries and the dual
costs. Consequently none of these changes evades the obstruction.

For example, at degree 35 the packet is \((1,8,15,22,29,30)\), and the
dual is \((1,1,1,1,1,2)/7\). No increase in exponent-search bounds can
produce a matrix of the earlier type with normalized weight sum greater
than 1.

\textbf{Interpretation.} This disproves universal applicability of the
specific positive-anticanonical matrix criterion proposed in v0.1.0. It
does not disprove the Hodge conjecture, exclude other uniruled carriers
or correspondences, or exclude non-anticanonical approaches.

\hypertarget{capacity_cycles-an-actual-cycle-for-the-missed-family}{%
\section{An actual cycle for the missed
family}\label{capacity_cycles-an-actual-cycle-for-the-missed-family}}

The construction now bypasses the carrier model entirely.

Fix \(q\ge2\), put \(m=5q\), and choose \(\kappa\in\mathbf C\) with

\[
\kappa^q=-1/5.
\]

Set

\[
z=x_5,\quad t_i=x_i^q\ (0\le i\le4),\quad P=x_0x_1x_2x_3x_4,
\]

and let \(e_j\) be the elementary symmetric polynomial of degree \(j\)
in \(t_0,\ldots,t_4\). Define

\begin{equation}\label{eq:capacity_cycles-4.1}
\boxed{
Z_\kappa=\{e_1=0,\ e_2=0,\ P-\kappa z^5=0\}\subset\mathbf P^5.
}\end{equation}

\hypertarget{capacity_cycles-proposition-4.1-containment-in-the-fermat-fourfold}{%
\subsection{Proposition 4.1: containment in the Fermat
fourfold}\label{capacity_cycles-proposition-4.1-containment-in-the-fermat-fourfold}}

Newton's identity is

\begin{equation}\label{eq:capacity_cycles-4.2}
p_5=e_1(e_1^4-5e_1^2e_2+5e_1e_3+5e_2^2-5e_4)-5e_2e_3+5e_5.\end{equation}

Since \(e_5=P^q\), on \(Z_\kappa\) one has

\[
F_m=p_5+z^{5q}=5P^q+z^{5q}
=(5\kappa^q+1)z^{5q}=0.
\]

Thus \(Z_\kappa\subset X_{5q}\). Its defining equations have degrees
\(q,2q,5\).

\hypertarget{capacity_cycles-proposition-4.2-the-surface-is-an-integral-complete-intersection}{%
\subsection{Proposition 4.2: the surface is an integral complete
intersection}\label{capacity_cycles-proposition-4.2-the-surface-is-an-integral-complete-intersection}}

Let

\[
S_q=\{e_1(x_0^q,\ldots,x_4^q)=e_2(x_0^q,\ldots,x_4^q)=0\}\subset\mathbf P^4.
\]

The coordinate-power map to the \(t\)-space is finite, so the dimension
is two. In characteristic zero, on \(e_1=0\) the equations are
equivalent to

\[
\sum x_i^q=0,\qquad\sum x_i^{2q}=0.
\]

At a hypothetical failure of Jacobian rank, there would be constants
\(\lambda,\mu\), not both zero, such that for each nonzero coordinate

\[
\lambda+2\mu x_i^q=0.
\]

The case \(\mu=0\) is impossible. Otherwise all nonzero \(x_i^q\) equal
the same nonzero constant. Their sum cannot vanish in characteristic
zero. Thus \(S_q\) is smooth. It is a connected complete intersection of
positive-dimensional ample divisors, hence irreducible.

The same argument with four coordinates shows that \(S_q\cap\{x_i=0\}\)
is a smooth reduced complete-intersection curve. Distinct such curves
share no irreducible component, since imposing two coordinate zeros
gives a zero-dimensional intersection. In particular the divisor of
\(P\) has valuation one at the generic point of each coordinate curve.

Forgetting \(z\) defines a finite flat cover \(Z_\kappa\to S_q\) of
degree five. On a chart trivializing \(\mathcal O_{S_q}(1)\), its
equation is \(z^5=P/\kappa\), with the corresponding trivializations
understood. At a generic coordinate curve its right side has valuation
one. The degree-five polynomial is Eisenstein over that discrete
valuation ring and is therefore irreducible over the function field of
\(S_q\). Finite flatness makes the cover generically integral without
vertical or embedded components, locally its monic algebra injects into
the generic field extension. Hence the cover is integral.

The two equations of \(S_q\), followed by the monic equation in \(z\),
form a homogeneous regular sequence. The resulting projective ideal is
saturated and defines the integral complete intersection \eqref{eq:capacity_cycles-4.1}.
Therefore

\[
\dim Z_\kappa=2,\qquad\deg Z_\kappa=q(2q)5=10q^2.
\]

The surface need not be smooth. Villaflor Loyola's complete-intersection
theorem does not impose smoothness of \(Z\), its proof explicitly
reduces to a smooth perturbation. We use this version, not an unverified
smooth-cycle variant.

\begin{figure}[tbp]\centering\includegraphics[width=0.98\linewidth]{30_explicit_fivephase_surface.pdf}\caption{A direct reconstruction of the known Aoki standard surface. Containment follows from exact Newton identities.}\end{figure}

\hypertarget{capacity_cycles-the-exact-jacobian-determinant}{%
\section{The exact Jacobian
determinant}\label{capacity_cycles-the-exact-jacobian-determinant}}

Define

\[
\begin{aligned}
f_1&=e_1,&g_1&=e_1^4-5e_1^2e_2+5e_1e_3+5e_2^2-5e_4,\\
f_2&=e_2,&g_2&=-5e_3,\\
f_3&=P-\kappa z^5,&g_3&=5\sum_{j=0}^{q-1}P^{q-1-j}\kappa^jz^{5j}.
\end{aligned}
\]

The geometric-series identity and \eqref{eq:capacity_cycles-4.2} give

\begin{equation}\label{eq:capacity_cycles-5.1}
F_m=f_1g_1+f_2g_2+f_3g_3.\end{equation}

Let

\[
H=(f_1,g_1,f_2,g_2,f_3,g_3),\qquad
D_\kappa=\det\operatorname{Jac}_{(x_0,\ldots,x_4,z)}H,
\]

and

\[
\Delta(t)=\prod_{0\le i<j\le4}(t_i-t_j).
\]

\hypertarget{capacity_cycles-theorem-5.1-uniform-determinant-formula}{%
\subsection{Theorem 5.1: uniform determinant
formula}\label{capacity_cycles-theorem-5.1-uniform-determinant-formula}}

\begin{equation}\label{eq:capacity_cycles-5.2}
\boxed{
D_\kappa=625q^5\Delta(x_0^q,\ldots,x_4^q)
\sum_{b=1}^{q-1}\kappa^{q-b}P^{b-1}z^{5(q-b)-1}.
}\end{equation}

\textbf{Proof.} First,

\begin{equation}\label{eq:capacity_cycles-5.3}
df_1\wedge dg_1\wedge df_2\wedge dg_2
=25\,de_1\wedge de_2\wedge de_3\wedge de_4.\end{equation}

Only the term \(-5de_4\) in \(dg_1\) survives after wedging with
\(de_1,de_2,de_3\), the row permutation has positive sign.

Next, direct differentiation yields

\begin{equation}\label{eq:capacity_cycles-5.4}
df_3\wedge dg_3=
25q\sum_{j=1}^{q-1}\kappa^jP^{q-1-j}z^{5j-1}\,dP\wedge dz.\end{equation}

For each interior index \(j\), the two contributing coefficients add to
\(j+(q-j)=q\). The endpoint terms have zero derivative coefficient. Thus
\eqref{eq:capacity_cycles-5.4} is an identity for every positive integer \(q\), not an
interpolation from checked values.

Finally,

\[
de_1\wedge\cdots\wedge de_5
=\Delta(t)\,dt_0\wedge\cdots\wedge dt_4.
\]

Use \(e_5=P^q\) and
\(dt_0\wedge\cdots\wedge dt_4=q^5P^{q-1}\,dx_0\wedge\cdots\wedge dx_4\).
Cancellation in the polynomial ring gives

\begin{equation}\label{eq:capacity_cycles-5.5}
de_1\wedge\cdots\wedge de_4\wedge dP
=q^4\Delta(x^q)\,dx_0\wedge\cdots\wedge dx_4.\end{equation}

Combining \eqref{eq:capacity_cycles-5.3}, \eqref{eq:capacity_cycles-5.4}, and \eqref{eq:capacity_cycles-5.5}, and putting \(b=q-j\), proves \eqref{eq:capacity_cycles-5.2}.
The identity extends across \(P=0\) because it is a polynomial identity.
\(\square\)

The degree of every monomial in \eqref{eq:capacity_cycles-5.2} is \(15q-6=3m-6\), as required for
a middle residue numerator.

\hypertarget{capacity_cycles-nonzero-periods-not-just-a-formal-character-match}{%
\section{Nonzero periods, not just a formal character
match}\label{capacity_cycles-nonzero-periods-not-just-a-formal-character-match}}

Fix \(1\le b<q\). In the order used in \eqref{eq:capacity_cycles-3.1}, the coefficient of

\begin{equation}\label{eq:capacity_cycles-6.1}
x_0^{b-1}x_1^{b+q-1}x_2^{b+2q-1}x_3^{b+3q-1}x_4^{b+4q-1}
 z^{5(q-b)-1}\end{equation}

in \(D_\kappa\) is

\begin{equation}\label{eq:capacity_cycles-6.2}
625q^5\kappa^{q-b}\ne0.\end{equation}

All permutations of the first five exponents have the same nonzero
magnitude, with the Vandermonde sign. Different values of \(b\) and
different permutations give different monomials. Every exponent lies
between zero and \(m-2\), so no term is killed by the Fermat Jacobian
ideal.

Define the dual numerator

\begin{equation}\label{eq:capacity_cycles-6.3}
Q_b=\prod_{i=0}^5x_i^{m-a_{q,b,i}-1}.\end{equation}

It has degree \(3m-6\). Because the product \(Q_bD_\kappa\) has socle
degree, a surviving monomial modulo \(J_{F_m}\) must have all six
exponents equal to \(m-2\). Formula \eqref{eq:capacity_cycles-6.1} is the unique term with this
complement. Therefore

\begin{equation}\label{eq:capacity_cycles-6.4}
Q_bD_\kappa\equiv625q^5\kappa^{q-b}S_m\pmod{J_{F_m}}.\end{equation}

For the Fermat polynomial,

\[
\det\operatorname{Hess}(F_m)=[m(m-1)]^6S_m.
\]

Apply Villaflor Loyola's Theorem 1.2 with dimension four and the
ordering fixed above. Its constant is

\[
c=\frac{625q^5\kappa^{q-b}}{[m(m-1)]^6}.
\]

The exact period is consequently

\begin{equation}\label{eq:capacity_cycles-6.5}
\boxed{
\int_{Z_\kappa}\operatorname{res}\left(\frac{Q_b\Omega}{F_m^3}\right)
=\frac{(2\pi i)^2}{50q}\,\kappa^{q-b}\ne0.
}\end{equation}

The normalization follows from \(625q^5/[2(5q)^6]=1/(50q)\). In
conventions using a differently normalized residue map, the displayed
constant changes with that convention, nonvanishing does not.

\hypertarget{capacity_cycles-rational-algebraicity-of-whole-packets}{%
\subsection{Rational algebraicity of whole
packets}\label{capacity_cycles-rational-algebraicity-of-whole-packets}}

Let \(h\) denote the hyperplane class. Since \(\deg Z_\kappa/m=2q\), the
primitive part

\[
\xi_\kappa=[Z_\kappa]-2q h^2
\]

is rational and algebraic. Villaflor Loyola's Theorem 1.1 identifies its
primitive residue numerator with \(-2D_\kappa/m\). Thus its complex
eigenspace support consists of exactly the \(120(q-1)\) characters
obtained from \eqref{eq:capacity_cycles-3.1} by permuting the first five entries and varying
\(b\).

That set is closed under multiplication by units modulo \(m\), by Lemma
3.1. Each of its eigenspaces has dimension one. The
\(\mathbf Q[G_m]\)-submodule generated by \(\xi_\kappa\) therefore has
complexification equal to the entire sum of those eigenlines: applying
the complex character idempotents isolates each nonzero component. Since
the generator and the group action are rational, this is the
complexification of the claimed rational submodule.

Every vector in that rational submodule is a rational linear combination
of classes of diagonal translates of \(Z_\kappa\) and of \(h^2\). Hence
every rational packet in the family is algebraic. Its rational algebraic
subspace has dimension

\[
120(q-1)
\]

for this fixed choice of the last coordinate. This is not the dimension
of all Hodge classes on \(X_{5q}\).

The conclusion recovers Aoki's standard family, including the elementary
decomposable midpoint case at even levels. It does not close arbitrary
nonstandard five-primary packets.

\hypertarget{capacity_cycles-a-concrete-example-at-degree-35}{%
\section{A concrete example at degree
35}\label{capacity_cycles-a-concrete-example-at-degree-35}}

Take \(q=7\), \(b=1\), \(\kappa^7=-1/5\). Then

\[
a=(1,8,15,22,29,30),\quad
Z_\kappa=\left\{
\begin{array}{l}
\sum_{i=0}^4x_i^7=0,\\
\sum_{0\le i<j\le4}x_i^7x_j^7=0,\\
x_0x_1x_2x_3x_4=\kappa x_5^5
\end{array}\right\}\subset X_{35}.
\]

The cycle has degree 490. The rational packet has dimension
\(\varphi(35)=24\), since its first character entry is 1. The
determinant coefficient is

\[
10,504,375\,\kappa^6.
\]

For

\[
Q_1=x_0^{33}x_1^{26}x_2^{19}x_3^{12}x_4^5x_5^4,
\]

one obtains

\[
\boxed{
\int_{Z_\kappa}\operatorname{res}\left(\frac{Q_1\Omega}{F_{35}^3}\right)
=\frac{(2\pi i)^2\kappa^6}{350}\ne0.
}
\]

At the same time, the uniform matrix-obstruction theorem above proves
that this packet has capacity exactly 1 and cannot use the
positive-anticanonical matrix route. Thus this is an explicit example
where failure of that constructor is compatible with, and bypassed by,
an actual algebraic cycle.

At degree 70 the standard packet \((1,15,29,43,57,65)\) obtained by
setting \(q=14\) is \textbf{different} from the earlier packet
\((1,20,24,42,61,62)\). Their capacities are respectively 1 and
\(11/10\). They must not be conflated.

\begin{figure}[tbp]\centering\includegraphics[width=0.98\linewidth]{31_socle_pairing35.pdf}\caption{At degree 35 the exact socle coefficient gives a nonzero period and a genuine packet component.}\end{figure}

\hypertarget{capacity_cycles-exact-computations-and-limitations}{%
\section{Exact computations and
limitations}\label{capacity_cycles-exact-computations-and-limitations}}

The final verifier uses integers, rational numbers, and symbolic
polynomial identities. It uses no floating-point arithmetic, random
sampling, network request, or optimization routine.

Its 137 grouped assertions include:

\begin{itemize}
\tightlist
\item
  balance and exact primal/global-dual certificates for seven specified
  packets,
\item
  the predecessor archive checksum,
\item
  Newton's identity and the elementary-symmetric Jacobian identity,
\item
  exact six-by-six intermediate determinants for \(q=2,3,4,5,7,11,17\),
\item
  independent socle-pairing checks for every \(b\) at \(q=2,7,11\),
\item
  all \(b\), all units, and all first-five permutations for
  \(2\le q\le64\), totaling 241,920 distinct support monomials and
  218,836 balance/conjugacy instances.
\end{itemize}

The seven globally certified capacity values are:

\begin{longtable}[]{@{}lll@{}}
\toprule\noalign{}
Modulus & Character & Capacity \\
\midrule\noalign{}
\endhead
\bottomrule\noalign{}
\endlastfoot
25 & (1,6,11,16,20,21) & 6/5 \\
33 & (1,4,16,22,25,31) & 2 \\
35 & (1,8,15,22,29,30) & 1 \\
39 & (1,7,16,22,34,37) & 2 \\
40 & (1,9,17,25,33,35) & 1 \\
55 & (1,12,23,34,45,50) & 1 \\
70 & (1,20,24,42,61,62) & 11/10 \\
\end{longtable}

These are seven specified inputs, not a census of all characters at
those levels. Capacities 2 in the table do not themselves supply a
six-column normal carrier.

The finite identities support the implementations. The all-parameter
assertions come from the separate mathematical proofs in the balance,
cycle and period arguments above, conditional only on the stated
standard algebraic-geometric and residue theorems. These proofs remain
at manuscript level.

\hypertarget{capacity_cycles-unclosed-targets}{%
\subsection{Unclosed targets}\label{capacity_cycles-unclosed-targets}}

No proof is given of the general Hodge conjecture, of coverage of every
Fermat packet, of the full CM half-twist Chow correspondence, or of the
earlier central-defect realization lemma. The universal
positive-anticanonical invariant-matrix conjecture is false by the
uniform matrix-obstruction theorem above and should not remain in a gap
ledger as an unproved plausible universal lemma.

The usable research architecture is now a collection of separately
certified constructors: carrier transfer where its hypotheses hold,
explicit standard surfaces where the capacity barrier is 1, and other
correspondences for packets not handled by either route. A proof that
such constructors exhaust all Hodge classes is still missing.
